\chapter*{Samenvatting}
\addcontentsline{toc}{chapter}{Samenvatting}

\glsunset{ADNEX}
\glsunset{RMI}
Klinische voorspelmodellen helpen artsen om in te schatten hoe groot de kans is dat iemand een bepaalde aandoening heeft. Dat is handig, maar alleen als zo'n model niet alleen gemiddeld goed werkt, maar ook voor nieuwe patiënten en in verschillende ziekenhuizen betrouwbare schattingen geeft.

Dit proefschrift gaat over vrouwen met een verdacht gezwel in de eierstokken of eromheen. Het centrale doel was om artsen beter te ondersteunen bij de vraag of zo'n gezwel waarschijnlijk goedaardig of kwaadaardig is, met bijzondere aandacht voor het \gls{ADNEX}-model dat vaak wordt gebruikt om het risico op eierstokkanker in te schatten.

Eerst onderzochten we hoe goed \gls{ADNEX} in andere studies heeft gewerkt. Een overzicht van 47 studies met meer dan 17.000 gezwellen liet zien dat het model meestal goed onderscheid maakt tussen goedaardige en kwaadaardige afwijkingen. Tegelijk bleek dat veel studies methodologisch niet ideaal waren opgezet en dat kalibratie, dus de vraag of de voorspelde kansen echt overeenkomen met de werkelijkheid, vaak niet goed werd nagekeken. Daarna vergeleken we \gls{ADNEX} rechtstreeks met de Risk of Malignancy Index (\gls{RMI}). Daarbij bleek \gls{ADNEX} duidelijk betrouwbaarder en nuttiger voor de klinische praktijk. Op basis van nieuwe multicentrische gegevens hebben we het model vervolgens aangepast tot \emph{ADNEX2}, zodat het beter past bij een tweestapsaanpak waarbij eerst gekeken wordt naar eenvoudige kenmerken die wijzen op een goedaardige massa.


Het proefschrift gaat ook over bredere vragen rond klinische voorspelmodellen. We onderzochten waarom sommige algoritmes op trainingsdata heel goed lijken, maar toch onbetrouwbare kansen geven wanneer ze echt worden ingezet. Ook ontwikkelden we methoden die rekening houden met verschillen tussen ziekenhuizen, zodat niet alleen een gemiddeld resultaat zichtbaar wordt, maar ook hoe sterk de prestaties van centrum tot centrum kunnen verschillen. Verder stelden we een raamwerk voor onzekerheid rond een individuele risicoschatting voor en ontwikkelden we een methode om alleen die variabelen te kiezen die echt nuttig zijn voor de zorg en niet onnodig veel kosten of metingen vragen.

De rode draad is eenvoudig: medische voorspelmodellen kunnen veel betekenen voor betere zorg, maar alleen als ze eerlijk worden getest, zorgvuldig worden gebruikt en hun onzekerheid duidelijk wordt uitgelegd. Dit proefschrift levert daarvoor nieuw bewijs, verbetert het bestaande \gls{ADNEX}-model en reikt methoden aan om toekomstige voorspelmodellen beter en veiliger te ontwikkelen.